\section{Conclusion}
\label{sec:conclusion}

We have closed the final gap in the framework for Verifiable
Institutions.

Paper~5~\cite{btv-cc} established that algorithmic systems can
embody the Separation of Powers as a structural invariant enforced
by linear types and cryptographic proofs. This paper establishes
that they can also embody the capacity for \emph{constitutional
evolution}---the one property that distinguishes a living
constitution from a fossil.

By formalising the \emph{Constitutional Amendment Protocol}, we
resolve the Capture Paradox: the system is rigid where rigidity
protects rights (Stone Clauses require Tripartite Ratification)
and flexible where flexibility enables justice (policy updates
require only Legislative authorisation). By encoding
\emph{Sunset Clauses} as linear resources, we invert the default
from legislative inertia to democratic renewal: authority must be
actively re-legitimised, not merely inherited.

The system now satisfies a new invariant, composing the amendment
mechanism into the type calculus of Papers 1--5:
\[
  \texttt{Change} \multimap
  (\texttt{Consensus} \otimes \texttt{Legitimacy})
\]
Evolution, like evidence, is a linear resource: it cannot be created
without consuming the consent of all three powers, and once consumed,
it produces a new constitutional state that is itself immutable and
permanently auditable.

The architecture of democracy is now complete. Paper~1 gave us the
obligation to prove. Paper~2 gave us the impossibility of forgetting.
Paper~3 gave us the right to privacy within accountability.
Paper~4 gave us the economics of honesty. Paper~5 gave us the
Separation of Powers. Paper~6 gives us the capacity to change---under
conditions that make change itself an act of governance.

From bits to Montesquieu to Jefferson: the machine that cannot be
corrupted, and the process by which it may, legitimately, be
rebuilt.
